In order to obtain estimates for the latent factors of the yield curve, we chose to apply principal component analysis (PCA), because it effectively distills the high-dimensional, highly correlated movements of the term structure into a set of independent latent factors. By decomposing the covariance matrix of yields, PCA recognizes that the curve behaves like a smooth, continuous function rather than a set of isolated data points. The mathematical requirement for orthogonality then forces these factors into the "Level, Slope, and Curvature" hierarchy. 

\subsection{Explained Variance}
Figure~\ref{fig:scree} presents the scree plot of the PCA applied to standardized yields across maturities. As expected, the first principal component explains approximately 97\% of total variance, while the second and third components explain roughly 2.5\% and 0.1\%, respectively. This indicates that the U.S. yield curve is overwhelmingly driven by a common level factor, although slope and curvature components remain economically meaningful despite their small variance shares.

\begin{figure}[H]
    \centering
    \includegraphics[width=0.7\textwidth]{figures/fig_PCA_scree.pdf}
    \caption{Explained variance ratio of yield curve principal components.}
    \label{fig:scree}
\end{figure}




\subsection{Factor Loadings}
Got get a better picture of the shape of the curve, we look at the factor loadings at all maturities used. Figure~\ref{fig:loadings} displays the PCA loading heatmap across maturities. The first principal component loads positively and nearly uniformly across all maturities, consistent with a level factor. The second component exhibits negative loadings at the short end and positive loadings at the long end, representing yield curve slope. The third component captures curvature, reflecting differential movements of the belly of the curve relative to short and long maturities.

\begin{figure}[H]
    \centering
    \includegraphics[width=0.85\textwidth]{figures/fig_PCA_loadings_heatmap.pdf}
    \caption{Heatmap of yield curve PCA loadings by maturity from 3 months until 30 years.}
    \label{fig:loadings}
\end{figure}


\subsection{Factor Dynamics and Recessions}

Figure~\ref{fig:pc-timeseries} plots the time series of the first three principal components with NBER recession shading. The level factor reflects long-run secular declines in interest rates and major monetary regime shifts. The slope factor exhibits pronounced swings around tightening cycles and yield curve inversions preceding recessions. The curvature factor shows episodic spikes during periods of financial stress and policy intervention.

\begin{figure}[H]
    \centering
    \includegraphics[width=0.9\textwidth]{figures/fig_PC_timeseries.pdf}
    \caption{Time series of yield curve principal components with recession shading.}
    \label{fig:pc-timeseries}
\end{figure}

\subsection{Rolling Loading Stability}

Figures~\ref{fig:roll1}--\ref{fig:roll3} present rolling-window PCA loading estimates for the first three components. While the overall factor structure is preserved, substantial drift is observed, particularly at the long end of the curve during the Global Financial Crisis and the COVID-19 period. These findings suggest that yield curve factors are regime-dependent and that long-maturity yields become structurally distorted during periods of unconventional monetary policy.

\begin{figure}[H]
    \centering
    \includegraphics[width=0.85\textwidth]{figures/fig_rolling_loading_drift_PC1.pdf}
    \caption{Rolling PCA loading drift for the level factor (PC1).}
    \label{fig:roll1}
\end{figure}

\begin{figure}[H]
    \centering
    \includegraphics[width=0.85\textwidth]{figures/fig_rolling_loading_drift_PC2.pdf}
    \caption{Rolling PCA loading drift for the slope factor (PC2).}
    \label{fig:roll2}
\end{figure}

\begin{figure}[H]
    \centering
    \includegraphics[width=0.85\textwidth]{figures/fig_rolling_loading_drift_PC3.pdf}
    \caption{Rolling PCA loading drift for the curvature factor (PC3).}
    \label{fig:roll3}
\end{figure}
